% Avsnitt 5.6, ur lectures/L05/README.md och appendix A.4.

\section{Sammanfattning}
\label{c5:sec:review}

\needspace{6\baselineskip}
\subsection*{Det här ska du kunna nu}
\begin{itemize}
  \item Säga vad som skiljer en \code{signal} från en \code{variable}, förutsäga det olika resultat
    var och en ger ur samma kod, och välja mellan dem medvetet.
  \item Förklara vad en FPGA fysiskt består av, varför två olika utseende beskrivningar av samma
    krets ger samma hårdvara, och varför ett Karnaughdiagram köper mindre här än på diskret logik.
  \item Förklara vad som begränsar en designs klockfrekvens, och varför fler register kan göra en
    design snabbare snarare än långsammare.
  \item Känna igen ett oavsiktligt lås på dess varning, säga vad som orsakade det och åtgärda det,
    och säga varför ”det kompilerade utan fel” lovar mindre för en FPGA-design än för ett
    C-program.
\end{itemize}

\needspace{6\baselineskip}
\subsection*{Frågor att testa dig själv med}
\begin{enumerate}
  \item En \code{signal} tilldelas två gånger, med olika värden, i ett varv genom en
    \code{process}. Vad hamnar på den, och när sker den uppdateringen egentligen?
  \item Samma fråga för en \code{variable}.
  \item Varför ger ackumulering över en loop med en \code{signal} ett annat, oftast felaktigt,
    resultat?
  \item Ingenting på en FPGA blir en grind. Vad blir en boolesk funktion i stället, och varför
    kostar en funktion av fyra variabler lika mycket oavsett om du minimerade den först?
  \item Vad är kritiska vägen, och varför kan fler vippor få en design att gå snabbare?
  \item Ett kombinatoriskt \code{case} tilldelar sin utgång i tre grenar av fyra. Vad bygger
    syntesverktyget för den fjärde, och varför är det det enda det \emph{kan} bygga?
  \item Varför är ett oavsiktligt lås en varning och inte ett fel, och varför gör det saken
    farligare snarare än mindre farlig?
\end{enumerate}

\needspace{6\baselineskip}
\subsection*{Blicken framåt}
Ingenting i \chapref{c6:ch} till \ref{c8:ch} inför något nytt idiom. Räknare, skiftregister och
timers håller alla tillstånd över klockflanker med hjälp av \chapref{c3:ch}:s klockade process och
schemaläggningsregeln ovan, och \chapref{c8:ch}:s tillståndsmaskiner lägger bara till en uppräknad
typ och ett \code{case} på den. Det som ändras är vad tillståndet används till.

Två saker att hålla utkik efter. \Chapref{c6:ch}:s \code{serial_rx8} är den första designen vars
korrekthet du inte kan kontrollera med ögat på en lysdiod, vilket är där de utdelade testbänkarna
slutar vara en formalitet. Och \secref{c4:sec:generics}:s generic sätts i arbete direkt:
\chapref{c7:ch}:s \code{timer} använder en för sitt tickantal, och \code{walking_led} instansierar
den timern vid sidan av \code{reset_sync} och \code{button_sync}, och skriver nästan ingen egen
logik.
